\thispagestyle{empty}


\noindent 
\textsf{\textbf{\Large Anotácia diplomovej práce}}

\medskip



\noindent
Slovenská technická univerzita v Bratislave \\
Fakulta elektrotechniky a informatiky

\bigskip
\noindent
\begin{tabular}{@{}l l}
Študijný odbor: & kybernetika \\
Študijný program: & Robotika a kybernetika \\
Autor: & Bc. Andrii Stoliar \\
Diplomová práca: & Návrh implementácie vybraných algoritmov riadenia \\
Školiteľ: & Ing. Marián Tárník, PhD. \\
Mesiac, rok odovzdania: & Máj, 2026 \\
\end{tabular}

\bigskip
\noindent
Kľúčové slová: riadenie systémov, laboratórny dynamický systém, identifikácia systému, numerická simulácia, PI regulátor, prepínanie pracovných bodov, gain scheduling, adaptívne riadenie, MRAC, gradientná metóda, MATLAB/Simulink, Arduino UNO, C++, PLC Siemens, TIA Portal, experimentálne overenie

\bigskip
\noindent
Táto diplomová práca sa zaoberá návrhom, implementáciou a experimentálnym
overením vybraných algoritmov riadenia pre laboratórny tepelný systém.
Hlavným cieľom práce je zostavenie riadiaceho systému pre reálny laboratórny
dynamický systém a realizácia vlastného softvérového návrhu na viacerých
implementačných platformách. V práci sú využité postupy z oblasti identifikácie
systémov, numerickej simulácie a vyhodnocovania kvality riadenia.

V prvej časti je opísaný použitý laboratórny systém, jeho vstupy, výstupy a spôsob
komunikácie s jednotlivými platformami. Na základe experimentálne nameraných
údajov je vykonaná identifikácia systému, pričom získaný model slúži ako podklad
pre návrh a simulačné overenie regulačných algoritmov. Následne sú navrhnuté
a realizované tri prístupy riadenia: PI regulátor, PI regulácia s prepínaním
pracovných bodov a adaptívna regulácia typu MRAC založená na gradientnej
metóde.

Jednotlivé algoritmy sú implementované v prostredí MATLAB/Simulink, v jazyku
C++ na platforme Arduino UNO a na priemyselnom PLC systéme Siemens
SIMATIC S7-1200 G2. Experimentálne overenie je zamerané na posúdenie
funkčnosti navrhnutých riadiacich systémov pri riadení reálneho zariadenia.
Výsledky práce ukazujú praktické možnosti implementácie zvolených algoritmov
na rôznych platformách a zároveň poukazujú na vplyv merania, časovania,
spracovania signálov a vlastností hardvéru na výsledné správanie regulačného
systému.


\pagebreak